\documentclass[a4paper,fleqn]{cas-sc}
\usepackage[utf8]{inputenc}
\usepackage[T5]{fontenc}
\usepackage[vietnamese]{babel}

\usepackage{amsmath}
\usepackage{amssymb}

\begin{document}
    \shorttitle{Biểu diễn SAT cho bài toán N-Queens}
    \shortauthors{Trần Việt Toàn}

    \title[mode=title]{Biểu diễn SAT tối ưu cho bài toán N-Queens: Đánh giá thực nghiệm các phương pháp mã hóa}

    \author[1]{Trần Việt Toàn}
    \address[1]{Trường Đại học Công nghệ - ĐHQGHN}

    \begin{abstract}
        Bài báo này trình bày phương pháp giải bài toán N-Queens bằng cách chuyển đổi thành bài toán thỏa mãn logic mệnh
        đề (SAT) thông qua thư viện PySAT. Trọng tâm của nghiên cứu là thiết kế và đánh giá các kỹ thuật mã hóa
        (encoding) để biểu diễn ràng buộc ``tối đa một'' (At-Most-One - AMO) và ``ít nhất một'' (At-Least-One - ALO).
        Chúng tôi triển khai phương pháp mã hóa Sequential Counter và Commander nhằm giảm thiểu số lượng mệnh đề sinh ra
        so với phương
        pháp Pairwise truyền thống. Kết quả của cấu hình SAT sẽ được thiết kế để đối chiếu thực nghiệm với các bộ giải
        chính xác (ILP, CP) như Cplex, Gurobi và CP-SAT trên các ma trận kích thước lớn.
    \end{abstract}

    \maketitle


    \section{Giới thiệu}\label{sec:intro}
    N-Queens là một bài toán tối ưu tổ hợp kinh điển, yêu cầu đặt $N$ quân Hậu lên bàn cờ $N \times N$ sao cho không
    có bất kỳ hai quân Hậu nào tấn công nhau. Để giải quyết bài toán này trên quy mô lớn ($N = 50, 100, 500$
    ), việc sử dụng
    các thuật toán quay lui (backtracking) thông thường là không khả thi. Thay vào đó, bài toán được dịch sang dạng
    chuẩn
    CNF (Conjunctive Normal Form) để giải bằng các công cụ SAT Solver hiện đại thông qua thư viện PySAT. Mục tiêu của
    báo
    cáo là so sánh hiệu năng của phương pháp SAT Solving với các phương pháp lập trình nguyên thủy (ILP) và lập trình
    ràng buộc (CP) thông qua các thông số cốt lõi: số lượng biến ($n_v$), số lượng mệnh đề ($n_c$), thời gian xây dựng
    mô hình và thời gian tìm kiếm.


    \section{Phương pháp biểu diễn (SAT Encoding)}\label{sec:sat_encoding}

    \subsection{Biến quyết định và Dạng chuẩn CNF}\label{subsec:variables}
    Gọi $X_{i,j}$ là biến Boolean đại diện cho ô ở hàng $i$, cột $j$ trên bàn cờ kích thước $N \times N$.
    \begin{equation}\label{eq:var_def}
        X_{i,j} =
        \begin{cases}
            1, & \text{nếu có quân Hậu đặt tại ô } (i,j) \\
            0, & \text{nếu ngược lại}
        \end{cases}
    \end{equation}

    Bài toán N-Queens được cấu thành từ hai loại ràng buộc cơ bản: ALO (At-Least-One - Ít nhất một) và AMO (At-Most-One
    - Tối đa một).

    \subsection{Ràng buộc ALO (At-Least-One)}\label{subsec:alo}
    Ràng buộc ALO yêu cầu mỗi hàng và mỗi cột phải có ít nhất một quân Hậu.
    Mã hóa cho một hàng $i$ có chiều dài $N$ được biểu diễn trực tiếp bằng một mệnh đề tuyển (clause) duy nhất:
    \begin{equation}\label{eq:alo_row}
        \bigvee_{j=1}^{N} X_{i,j}
    \end{equation}

    \subsection{Tối ưu Ràng buộc AMO bằng Sequential Counter}\label{subsec:amo_seq}
    Ràng buộc AMO cấm việc đặt nhiều hơn một quân Hậu trên cùng một hàng, cột hoặc đường chéo.
    Nếu sử dụng phương pháp Binomial (Pairwise) cơ bản, việc bắt cặp tất cả các biến sẽ sinh ra $\mathcal{O}(N^2)$
    mệnh đề, gây bùng nổ bộ nhớ khi $N$ lớn.

    Để tối ưu, chúng tôi cài đặt phương pháp \textbf{Sequential Counter} (Bộ đếm tuần tự).
    Phương pháp này giới thiệu $N-1$ biến phụ (auxiliary variables) $S_1, S_2, \dots, S_{N-1}$ cho mỗi dãy $N$ ô.
    Biến $S_i$ mang ý nghĩa trạng thái: $S_i = \text{True}$
    khi và chỉ khi tồn tại ít nhất một biến mang giá trị True trong tập tiền tố $(X_1, X_2, \dots, X_i)$.

    Hệ thống luật của Sequential Counter được dịch sang CNF như sau:
    \begin{itemize}
        \item \textbf{Luật biên trái:} Nếu ô đầu tiên có Hậu, tín hiệu truyền đi được kích hoạt.
        \begin{equation}\label{eq:seq_left}
            X_1 \to S_1 \equiv \neg X_1 \lor S_1
        \end{equation}

        \item \textbf{Luật truyền tin (với $1 < i < N$):} Trạng thái được cập nhật và truyền tiếp dọc theo chuỗi.
        \begin{equation}\label{eq:seq_mid_prop}
            (X_i \lor S_{i-1}) \to S_i \equiv (\neg X_i \lor S_i) \land (\neg S_{i-1} \lor S_i)
        \end{equation}

        \item \textbf{Luật cấm AMO (với $1 < i < N$):} Nếu tín hiệu từ tiền tố báo đã có Hậu, cấm đặt Hậu ở ô hiện tại.
        \begin{equation}\label{eq:seq_mid_ban}
            S_{i-1} \to \neg X_i \equiv \neg S_{i-1} \lor \neg X_i
        \end{equation}

        \item \textbf{Luật biên phải:} Áp dụng lệnh cấm cho ô cuối cùng.
        \begin{equation}\label{eq:seq_right}
            S_{N-1} \to \neg X_N \equiv \neg S_{N-1} \lor \neg X_N
        \end{equation}
    \end{itemize}

    \subsection{Tối ưu Ràng buộc AMO bằng Commander Encoding}\label{subsec:amo_cmd}
    Khi quy mô bàn cờ tăng lên, phương pháp Sequential có thể tạo ra chuỗi lan truyền dài. Để kiểm soát cấu trúc biến
    phụ, chúng tôi áp dụng phương pháp \textbf{Commander Encoding}. Kỹ thuật này chia các biến thành nhóm, tạo một biến
    đại diện cho mỗi nhóm rồi áp đặt AMO lên các đại diện đó; phương pháp này giúp tạo tính cục bộ (locality) theo
    khối và cực kỳ phù hợp cho các miền biến vừa và lớn[cite: 13]. Cho tập hợp $n$ biến, thuật toán phân hoạch chúng
    thành các nhóm $G_1, G_2, \dots, G_m$ có kích thước cố định (ví dụ$g = 3$). Mỗi nhóm $G_i$ được gán một biến chỉ huy
    (commander variable) $C_i$. Hệ thống luật được quy đổi sang CNF bao gồm:
    \begin{itemize}
        \item \textbf{Luật Chỉ huy Sai:}
        Nếu biến chỉ huy mang giá trị False, tất cả biến trong nhóm phải mang giá trị False.
        \begin{equation}
            C_i \lor \neg X_k \quad (\forall X_k \in G_i)
        \end{equation}
        \item \textbf{Luật Chỉ huy Đúng:}
        Nếu biến chỉ huy mang giá trị True, có chính xác một biến trong nhóm mang giá trị True. Điều kiện này bao gồm
        ALO cục bộ:
        \begin{equation}
            \neg C_i \lor \bigvee_{X_k \in G_i} X_k
        \end{equation}
        và AMO nội bộ bằng cách áp dụng phương pháp Pairwise cho các cặp biến $(X_a, X_b)$ bên trong $G_i$:
        $\neg X_a \lor \neg X_b$.
        \item \textbf{Ràng buộc trên Chỉ huy:} Bản thân các biến chỉ huy $C_1, \dots, C_m$
        tiếp tục phải tuân thủ điều kiện AMO. Nếu số lượng chỉ huy lớn, quá trình gom nhóm được áp dụng đệ quy.
    \end{itemize}

\end{document}